\section*{第 \ref{ch:b2:dispersion-wave-packets} 章——色散与波包}

\begin{solution}{exo:b2:dispersion-wave-packets:1}
(a) $k = \qty{0.063}{rad/m}$, $\omega = \sqrt{gk} = \qty{0.79}{rad/s}$, $T = \qty{8.0}{s}$;
$v_\varphi = \qty{12.5}{m/s}$, $v_g = \qty{6.3}{m/s}$. (b) $\omega = \qty{0.45}{rad/s}$, $k = \omega^2
/g = \qty{0.021}{rad/m}$, $\lambda = \qty{310}{m}$; $v_\varphi = gT/2\pi = \qty{22}{m/s}$, $v_g =
\qty{11}{m/s}$；$5 \times 10^6/11 = \qty{4.6e5}{s} = 5.3$ 天。(c) $\lambda = \qty{6}{m}$，$v_g
= \qty{1.6}{m/s}$：需耗时五周，而短波早在此之前
就已被黏性\omterm{def:b2:dispersion-wave-packets:complex}{衰减}，并被风打散。
\end{solution}

\begin{solution}{exo:b2:dispersion-wave-packets:2}
$\omega = \qty{1.26e8}{rad/s}$, $\omega_c = \qty{6.3e7}{rad/s}$: $k = \sqrt{\omega^2 - \omega_c^2}/c =
\qty{0.36}{rad/m}$, $v_\varphi = \qty{3.5e8}{m/s}$, $v_g = c^2k/\omega = \qty{2.6e8}{m/s}$,
二者之积为 $c^2$。在 \qty{5}{MHz} 时，$\omega < \omega_c$：$\kappa = \sqrt{\omega_c^2 - \omega^2}/c =
\qty{0.18}{m^{-1}}$，深度为 \qty{5.5}{m}。
\end{solution}

\begin{solution}{exo:b2:dispersion-wave-packets:3}
拍频为 \qty{4}{Hz}，响度周期为 \qty{0.25}{s}。水波：$f = 0.125$ 和
\qty{0.119}{Hz}；包络为 $\cos(\delta\omega t - \delta kx)$，取 $\delta k = \qty{2.9e-3}{rad/m}$
（相邻拍相距 \qty{1.1}{km}），速率为 $\delta\omega/\delta k = \qty{6.4}{m/s}$——
即 \qty{105}{m} 处的\omterm{prop:b2:dispersion-wave-packets:group}{群速度}，$\tfrac12\sqrt{g\lambda/2\pi} = \qty{6.4}{m/s}$。
\end{solution}

\begin{solution}{exo:b2:dispersion-wave-packets:4}
$\ln(b/a) = 1.20$: $\Lambda = \qty{0.24}{\micro H/m}$, $\Gamma = \qty{106}{pF/m}$, $c = 1/
\sqrt{\Lambda\Gamma} = \qty{2.0e8}{m/s}$, $Z_c = \qty{48}{\ohm}$, \qty{5.0}{ns/m}; \qty{10}{ns}
$\leftrightarrow$ \qty{2.0}{m}。空气中：$\Gamma = \qty{46}{pF/m}$，$c = c_0$，$Z_c = \qty{72}{\ohm}$。
\end{solution}

\begin{solution}{exo:b2:dispersion-wave-packets:5}
(a) $\omega^2 = c^2k^2 + (EI/\mu)k^4$; $v_\varphi = \sqrt{c^2 + (EI/\mu)k^2}$; $v_g = (c^2k +
2(EI/\mu)k^3)/\omega$. (b) $c^2 = \qty{1.12e5}{m^2/s^2}$, $k = \qty{66}{rad/m}$, $(EI/\mu)k^2
= \qty{7.0e3}{m^2/s^2}$：$v_\varphi/c = \sqrt{1.063} = 1.031$，$+3.1\%$——即
\cref{ch:b2:waves-on-strings} 中的 $53$ 音分。(c) $v_g > v_\varphi$：反常的；高次
谐波跑在前面，拨弦的起音转为一段短促的
下降啁啾。
\end{solution}

\begin{solution}{exo:b2:dispersion-wave-packets:6}
(a) $-\mu\omega^2 + \iu\alpha\omega = -Tk^2$: $\underline k^2 = (\omega^2/c^2)(1 - \iu\alpha/\mu\omega)$.
(b) $\underline k \approx (\omega/c)(1 - \iu\alpha/2\mu\omega)$: $k' = \omega/c$, $k'' = \alpha/2\mu c =
\alpha/2Z$; $8.7k''$ dB/m. (c) $c = \qty{100}{m/s}$, $Z = \qty{0.5}{kg/s}$, $\mu\omega = 3.1 \gg
0.02$：$k'' = \qty{0.02}{m^{-1}}$，$\delta = \qty{50}{m}$，\qty{0.17}{dB/m}。(d) 此时不属于该
极限（$k' = \omega/c$）；当 $\alpha \gtrsim \mu\omega$ 时，平方根不再
线性化，且 $k'$ 对 $\omega$ 呈非线性依赖：色散的（即
电缆的扩散极限）。
\end{solution}

\begin{solution}{exo:b2:dispersion-wave-packets:7}
(a) $v_\varphi = \sqrt{\gamma k/\rho}$，$v_g = \tfrac32v_\varphi$：群超过
波峰。(b) $v_\varphi^2 = g/k + \gamma k/\rho$ 在 $k_m = \sqrt{\rho g/\gamma}$ 处取极小值：
$\lambda_m = 2\pi\sqrt{\gamma/\rho g} = \qty{1.7}{cm}$, $v_{\min} = (4g\gamma/\rho)^{1/4} =
\qty{23}{cm/s}$。(c) 尾迹是相对于
船静止的图样，由\omterm{def:b2:dispersion-wave-packets:relation}{相速度}与船速相同的波构成；其中没有一个
慢于 \qty{23}{cm/s}。(d) 毛细区，$v_g > v_\varphi$：波峰在
环的前沿生成，在后沿消失。
\end{solution}

\begin{solution}{exo:b2:dispersion-wave-packets:8}
(a) $k = \qty{0.126}{rad/m}$, $\omega'' = -\tfrac14\sqrt{g/k^3} = \qty{-18}{m^2/s}$: $t_{\text{展宽}}
= 2.5 \times 10^5/18 = \qty{1.4e4}{s}$，四小时。(b) \qty{10}{ps}：$t_{\text{展宽}} =
4 \times 10^{-6}/0.16 = \qty{25}{\micro s}$，\qty{5}{km}，取 \qty{2e8}{m/s}；\qty{1}{ns}：$0.04/0.16 =
\qty{0.25}{s}$，\qty{50000}{km}。短脉冲（高比特率）在
数公里之内便展宽：长距离链路使用补偿光纤，或采用中等速率。
\end{solution}

\begin{solution}{exo:b2:dispersion-wave-packets:9}
(a) $\partial_x^2v = R\Gamma\,\partial_tv$, $D = 1/R\Gamma = \qty{1.1e12}{m^2/s}$. (b) $\underline k^2
= -\iu R\Gamma\omega$: $k' = k'' = \sqrt{R\Gamma\omega/2}$; \qty{1}{Hz}: $k = \qty{1.7e-6}{m^{-1}}$,
$v_\varphi = \qty{3.7e6}{m/s}$, $\delta = \qty{600}{km}$; \qty{10}{Hz}: $v_\varphi = \qty{1.2e7}{m/s}$,
$\delta = \qty{190}{km}$. (c) $R\Gamma L^2 = 9 \times 10^{-13} \times 9 \times 10^{12} = \qty{8}{s}$
每点：约每分钟一个字。(d) $\Lambda = R\Gamma/G = \qty{9}{H/km}$——一个
荒谬的电感（对海底电缆而言）；加感线圈曾用于陆地
电话线路，其 $G$ 与 $R\Gamma$ 不同。
\end{solution}

\begin{solution}{exo:b2:dispersion-wave-packets:10}
(a) $\partial_te = \dot y(\mu\ddot y + Ky) + Ty'\dot y' = \dot yTy'' + Ty'\dot y' = \partial_x
(Ty'\dot y) = -\partial_x\mathcal P$. (b) $\langle e\rangle = \tfrac14A^2(\mu\omega^2 + Tk^2 + K) =
\tfrac12\mu\omega^2A^2$（利用 $\mu\omega^2 = Tk^2 + K$）；$\langle\mathcal P\rangle = \tfrac12Tk\omega
A^2$。(c) 比值 $Tk/\mu\omega = c^2k/\omega = v_g$。(d) $y = A\eu^{-\kappa x}\cos\omega t$：
$\mathcal P \propto \sin2\omega t$，平均值为零——没有能量穿过介质。
\end{solution}

\begin{solution}{exo:b2:dispersion-wave-packets:11}
(a) $\tanh kh \to 1$：$\omega^2 = gk$；$kh \ll 1$：$\omega = \sqrt{gh}\,k$，速率为 $\sqrt{gh}$，
非色散的。(b) $kh = 0.13$：浅水；$c = \sqrt{9.81 \times 4000} = \qty{198}{m/s}$，
$T = \qty{17}{min}$；\qty{8000}{km}，用时 \qty{11}{h}。(c) $\tfrac12\rho gA^2 = \qty{1.2}{kJ/m^2}$
在 $200 \times 1000$ km$^2$ 范围内：\qty{2.5e14}{J}。格林：$A \propto h^{-1/4}$，$(4000/10)^{1/4}
= 4.5$：\qty{2.2}{m}，陡化与爬高之前。(d) 在 $h \approx \lambda/2 =
\qty{50}{m}$；位于较浅水域中的波峰部分减速，因此
波峰转向，与等深线平行——即折射。
\end{solution}

\begin{solution}{exo:b2:dispersion-wave-packets:12}
(a) 在 $x = \ell$ 处：$v = v_i + v_r$，$i = (v_i - v_r)/Z_c$，$v = R_Li$：$r = (R_L -
Z_c)/(R_L + Z_c)$。(b) $0$，$-1$，$+1$，$0.5$。(c) 线路的输入端看上去像
$Z_c$，因此示波器在 $t = 0$ 时读到 \qty{0.5}{V}，随后回波 $0.5r$ V 在
$2\ell/c = \qty{1}{\micro s}$ 时刻分别为：无、\qty{-0.5}{V}、\qty{+0.5}{V}、\qty{+0.25}{V}；
匹配的信号源吸收了回波，因此不会有第二个回波。(d)
$\tfrac12 \times 2 \times 10^8 \times 1.8 \times 10^{-6} = \qty{180}{m}$.
\end{solution}

\begin{solution}{pb:b2:dispersion-wave-packets:1}
\textbf{1.} $\partial_xv = -\Lambda\partial_ti - Ri$，$\partial_xi = -\Gamma\partial_tv$；设
$\eu^{\iu(\omega t - kx)}$：$\iu kv = (\iu\omega\Lambda + R)i$，$\iu ki = \iu\omega\Gamma v$，由此
$k^2 = \omega\Gamma(\omega\Lambda - \iu R)$.

\textbf{2.} $\omega = R/\Lambda = \qty{4000}{rad/s}$，\qty{640}{Hz}：在几赫兹时 $\Lambda\omega
\ll R$，属于扩散型。

\textbf{3.} $k' = k'' = \sqrt{R\Gamma\omega/2}$，其中 $R\Gamma = \qty{5e-13}{s/m^2}$：\qty{1}{Hz}：
$\qty{1.25e-6}{m^{-1}}$, $v_\varphi = \qty{5e6}{m/s}$, $\delta = \qty{800}{km}$; \qty{4}{Hz}:
$\qty{2.5e-6}{m^{-1}}$, $\qty{1e7}{m/s}$, \qty{400}{km}.

\textbf{4.} $8.7L/\delta$：\qty{38}{dB}（\qty{1}{Hz}），\qty{76}{dB}（\qty{4}{Hz}）：只有
最慢的分量得以保留；尖锐的“点”到达时变成长而低的驼峰。

\textbf{5.} $R\Gamma L^2 = 5 \times 10^{-13} \times 1.2 \times 10^{13} = \qty{6}{s}$：每隔
六秒一个“点”，约每分钟一个词（真实的电缆
配有灵敏的镜式检流计，只能做到几个词）。

\textbf{6.} $\Lambda = R/\omega = \qty{80}{mH/km}$。则 $k^2 = \Gamma R\omega(1 - \iu)$，$k = \sqrt{\Gamma R\omega}
\,2^{1/4}\eu^{-\iu\pi/8}$: $k' = \qty{3.9e-6}{m^{-1}}$, $k'' = \qty{1.6e-6}{m^{-1}}$: $v_\varphi =
\qty{6.4e6}{m/s}$，$\delta = \qty{620}{km}$，整条电缆上为 \qty{49}{dB}，而非
\qty{76}{dB}.

\textbf{7.} 每隔 \qty{6}{s} 一个“点”约合 \qty{0.2}{bit/s}：光纤快
$10^{10}$ 至 $10^{11}$ 倍。

\textbf{8.} $v_\varphi = gT/2\pi$, $v_g = gT/4\pi$.

\textbf{9.} \qty{18}{s}: $v_g = \qty{14.1}{m/s}$, $4 \times 10^6/14.1 = \qty{3.3}{\text{天}}$;
\qty{6}{s}：\qty{4.7}{m/s}，\qty{9.9}{\text{天}}：涌浪持续时间超过六天。

\textbf{10.} $t = D/v_g = 4\pi D/gT$，于是 $1/T = gt/4\pi D$。$\Delta(1/T) = 1/12 - 1/16
= \qty{0.0208}{Hz}$，在 \qty{86400}{s} 内：$g/4\pi D = \qty{2.4e-7}{s^{-2}}$，$D = \qty{3200}{km}$。

\textbf{11.} $\tfrac12\rho gA^2 = \qty{11}{kJ/m^2}$；能量以 $v_g = gT/4\pi =
\qty{9.4}{m/s}$ 传播，每米波峰输送 \qty{103}{kW}。

\textbf{12.} $10^3 \times 103 \times 0.3 = \qty{31}{MW}$：相当于十台风力涡轮机。

\textbf{13.} $\lambda = 2\pi U^2/g = \qty{16}{m}$。每个波分量的能量
以其\omterm{def:b2:dispersion-wave-packets:relation}{相速度}的一半传播，图样因此限于船后的一个楔形区域内
（半角 \qty{19.5}{\degree}，与船速无关）。

\textbf{14.} 一个波峰在波群后端产生，赶上它
（波峰相对于水的速度为 $2v_g$，相对于波群的速度为 $v_g$
）并在前端消亡。波群长 $120 \times 9.4 = \qty{1.1}{km}$，
约为五个波长（$\lambda = gT^2/2\pi = \qty{225}{m}$）；十个波峰经过
浮标，每个波峰存在约两分钟。

\textbf{15.} $kh = 0.13$：浅水。$c = \qty{198}{m/s}$，$T = \lambda/c = \qty{17}{min}$，即
\qty{8.4}{h}.

\textbf{16.} 非色散的：所有分量均以 \qty{198}{m/s} 速率传播，
一起作为一个脉冲到达（涌浪速率约为 $\sim\qty{10}{m/s}$，若按周期分开，
则需要一周）。

\textbf{17.} $\tfrac12\rho gA^2 = \qty{1.8}{kJ/m^2}$，在 $500 \times 200$ km$^2$ 范围内总能量为 \qty{1.8e14}{J}，
相当于一兆吨的 $4\%$。

\textbf{18.} $A = 0.6(4000/20)^{1/4} = \qty{2.3}{m}$; $c = \sqrt{gh} = \qty{14}{m/s}$, $\lambda =
cT = \qty{14}{km}$。波峰所在处的水比其前方的波谷处更深，
传播更快并追上它：波前陡化，形成涌潮。

\textbf{19.} 水面在一刻钟内上升 \qty{0.6}{m}，在坡度为
$10^{-5}$ 的斜坡上，人毫无感觉。

\textbf{20.} 平均速率为 $\sqrt{gh}$（取平均水深）；由于
$\sqrt h$ 为凹函数，浅水段使波减慢的程度大于深水段使其加快的程度，
故路径平均速率小于 $\sqrt{g\bar h}$。

\textbf{21.} $\omega' = \tfrac12\sqrt{g/k}$, $\omega'' = -\tfrac14\sqrt{g/k^3}$; $k = \qty{0.063}{rad/m}$:
$\omega'' = \qty{-50}{m^2/s}$.

\textbf{22.} $t_{\text{展宽}} = 10^6/50 = \qty{2e4}{s}$，\qty{5.6}{h}；波包
已移动 $v_gt = 6.25 \times 2 \times 10^4 = \qty{125}{km}$。

\textbf{23.} 一段长的波列包含一个窄的 $k$ 带，其\omterm{prop:b2:dispersion-wave-packets:group}{群速度}
几乎重合：各分量保持在一起的时间更长。
海啸的介质为非色散的：速率相等，不发生展宽。

\textbf{24.} $\omega'' = \hbar/m = \qty{1.2e-4}{m^2/s}$：$t_{\text{展宽}} = 10^{-18}/1.2 \times
10^{-4} = \qty{9e-15}{s}$; $v_g = \hbar k/m = \qty{7.3e6}{m/s}$: \qty{60}{nm}.

\textbf{25.} 电缆：$k^2 = -\iu R\Gamma\omega$，$v_\varphi = \sqrt{2\omega/R\Gamma}$，无
明确的 $v_g$，\omterm{def:b2:dispersion-wave-packets:complex}{衰减}并弥散。涌浪：$\omega^2 = gk$，$v_\varphi = 2v_g$，
展宽。海啸：$\omega = \sqrt{gh}\,k$，$v_\varphi = v_g$，均不展宽。电子：
$\omega = \hbar k^2/2m$，$v_g = 2v_\varphi$，展宽。
\end{solution}
